\section{Supplementary Protocol}

\subsection{Corpus and leakage controls}
\label{sec:corpora_controls}

Dreaddit, GoEmotions, CaChe, and ParlaMint-GB enter the study protocol. The Dreaddit and GoEmotions official training splits support development, their test splits are reserved for in-domain audits, and all 11 CaChe focus groups and the ParlaMint-GB debates are reserved for held-out cross-domain evaluation. The CaChe comparison remains the sole confirmatory test, and the GoEmotions and ParlaMint-GB results are descriptive. ``Held out'' describes separation within this study, not model-pretraining exposure. Dreaddit and GoEmotions labels and published CaChe themes are neither scoring gold nor routing features.

Table~\ref{tab:adopted_corpora} reports deterministic preparation inventories rather than approved samples or analysis denominators. Dreaddit exclusions preserve post clusters; because the release has no author identifier, the study makes no participant-level claim. CaChe sampling preserves focus-group and transcript-local speaker clusters; eligible participant turns supply candidate content, while short and collective turns serve only as context. GoEmotions exclusions preserve comment clusters, since the simplified split carries no thread linkage. ParlaMint-GB sampling preserves debate and speaker clusters, and speaker turns supply context.

Model processing, rater display, or packet generation remains blocked until corpus-specific institutional, source or platform, provider, access, retention, and privacy requirements are documented. Prompts, packet rules, routing features, thresholds, and actor manifests must be frozen before held-out generation. Every excerpt and context window shown to a rater or considered for release requires two-person privacy and permissions review.

\subsection{Reviewers, assignment, and freezes}
\label{sec:reviewer_protocol}

RQ1 distinguishes researchers without specialist expertise, qualitative-methods experts, and domain experts. Eligibility, handling of dual expertise, recruitment, compensation, exclusions, and independence safeguards are fixed before recruitment. Reviewers receive the same tutorial, blinded packet, and response form. Assignment is balanced by corpus and flaw family, gives a reviewer at most one version of a base packet, and keeps methods and domain ratings independent until adjudication.

RQ2 uses prospectively frozen prompts corresponding to researchers without specialist expertise, qualitative methods experts, and domain experts, with three separately identified reviewer snapshots from different model families. One snapshot is designated primary, and the other two are sensitivity reviewers. Candidate-generation and reviewer models have distinct actor and run records. Reviewer payloads exclude defect labels, human responses, adjudication, router decisions, and outcomes, and all LLM responses are sealed before human ratings. The preregistration fixes prompts, decoding, repetitions, aggregation, retries, terminal-failure handling, exclusions, assignment, and the rule for selecting role outputs. \fillin{INSERT FINAL RECRUITMENT, ASSIGNMENT, ACTOR, PROMPT, AND OUTPUT-SEAL IDENTIFIERS FROM THE FROZEN RECORD}.

RQ3 prospectively freezes the eligible defective packets, revision model and prompt, feedback-selection and formatting rules, conditions, allocation, failure handling, repair-panel eligibility, blinding, outcome rule, and analysis procedure. The no-feedback and routed-feedback conditions use the same starting packet and revision process. A no-specialist route uses the initial researcher's locked feedback, a one-specialist route uses that specialist's locked feedback, and a both-specialists route presents both locked records in a fixed order. The repair panel sees neither the feedback source nor the condition and cannot evaluate feedback it produced. \fillin{INSERT FINAL RQ3 REVISION, FEEDBACK, ALLOCATION, PANEL, AND OUTPUT-SEAL IDENTIFIERS FROM THE FROZEN RECORD}.

\subsection{Statistical implementation}
\label{sec:artifact_metrics}

Section~\ref{sec:evaluation_methods} defines the metrics, diagnostic measures, estimands, and result classification. Section~\ref{sec:evaluation_analysis} defines scoring, miss rules, comparison design, and sensitivity status. The procedures below specify route application, resampling, multiplicity, and conditional inference.

WarrantRoute applies the frozen route offline to the RQ2 LLM outputs. A route with no specialist selects the output from the researchers prompt. Methods and domain routes select the corresponding role output. A route to both takes the union of methods and domain flags. A failed selected role contributes no flag and is never replaced by another role. The fixed role selected on development data is not reselected on held out data or inside a bootstrap resample.

The paired bootstrap resamples the source connected component fixed in the packet manifest. All linked base packets, controlled variants, role outputs, model snapshots, routes, and repetitions travel together. Dreaddit resampling preserves post and document linkage. CaChe resampling preserves focus group and transcript local speaker dependence. GoEmotions resampling preserves comment-level linkage. ParlaMint-GB resampling preserves debate and speaker linkage. If a stratum has too few independent clusters or yields a degenerate bootstrap distribution, the manuscript reports \(TP/N\) and ``CI not estimable'' rather than substituting an unclustered binomial interval.

The preregistration defines each confirmatory comparison family and its multiplicity rule before outcome access. Any additional model, error type, feedback source, or repair condition comparisons remain descriptive unless they are prospectively included in a closed family. The stability sensitivity retains a role flag selected in at least two of three responses, reapplies the same route, and reports its paired change from repetition 1. \fillin{INSERT FINAL SAMPLE SIZE BASIS, SOURCE CLUSTER DEFINITIONS, PRIMARY MODEL SNAPSHOT, FIXED ROLE, EXCLUSIONS, AND ANALYSIS EXPORT IDENTIFIER}.

After RQ1, the frozen RQ1 WarrantRoute policy supplies the route for each common packet. For both arms, a route with no specialist selects the output from researchers without specialist expertise or from the corresponding prompt for the designated primary LLM reviewer. Methods and domain routes select the corresponding output. A route to both takes the union of methods and domain flags. The common packet manifest designates one primary human rating per packet and role before outcomes. Additional human ratings enter reviewer sensitivity analyses. The paired interval is conditional on the fixed recruited reviewers and does not generalize to a reviewer population.

RQ3 retains every allocated defective packet, including packets on which the routed review misses the verified error. All revision conditions for a base packet travel together when source clusters are resampled. A terminal revision or repair-assessment failure counts as an unsuccessful repair and is not dropped or replaced. Target repair, preservation of accurate material, and collateral errors are also reported separately so that the composite outcome cannot conceal the reason for failure.

The result classification follows Section~\ref{sec:evaluation_methods}. Every interpretation is limited to the fixed packets, recruited reviewers, frozen models and prompts, and prespecified detection or repair outcome.

\section{Data Governance}
\label{sec:detailed_ethics}

Approval of an original collection or a public license does not authorize this secondary study, external model processing, rater disclosure, quotation, retention, or release. No adopted corpus is described as anonymous. Dreaddit and GoEmotions contain searchable verbatim public text from social media. CaChe contains selectively redacted, sensitive sexual and reproductive health narratives. ParlaMint-GB contains public parliamentary speech attributable to named officials. Restricted copies use project-pseudonymous identifiers, minimize fields, preserve protected linkage data separately, and retain only the context required for source-grounded review and repair assessment. Protected text, feedback, and revised outputs are sent only to services whose processing, retention, and training terms are approved for the corpus and purpose.

No excerpt is shown to reviewers or repair assessors or released without two-person privacy and permissions review. Raters participate voluntarily under the approved consent and withdrawal procedure, outside supervisory, employment, grading, or authorship decisions. Public artifacts are limited to aggregate results, code, prompts, schemas, non-text manifest identifiers, and material specifically cleared for release. Routing and repair are research-review aids, not clinical or participant-level decision systems. \fillin{INSERT THE FINAL INSTITUTIONAL DETERMINATION, APPROVED PROVIDERS, ACCESS ROLES, RETENTION AND DELETION RULES, RECRUITMENT TERMS, AND RELEASE CONDITIONS}.

\section{Supplementary Reporting}
\label{sec:appendix_results}

No experimental result is available. Authorized analyses report reconciled recall and repair denominators with frozen actor, prompt, packet, revision, and analysis identifiers. Tables~\ref{tab:rq2_accounting} and~\ref{tab:rq2_role_flaw} define the RQ2 result organization, and the RQ3 reporting block below defines the repair accounting. Every field remains a placeholder until an authorized frozen real-corpus export supplies it.

\ifpredatadraft
\subsection{Working ParlaMint-GB interim check}

The following working tables are for local pipeline inspection only. They are not manuscript results, were not produced from a frozen authorized analysis export, and must not be interpreted as confirmatory evidence. The packet bank samples ParlaMint-GB from the full local speaker-turn population: 670,912 speaker turns, 642,742 turns with at least eight words, and 607,907 prompt-usable turns after the local reviewer context filter of at most 800 words. The displayed run uses Qwen3 8B on deterministic working packets with five flaw families balanced by construction.

\begin{table*}[t]
\centering
\scriptsize
\setlength{\tabcolsep}{3.2pt}
\renewcommand{\arraystretch}{1.05}
\begin{tabular}{@{}lccc@{}}
\toprule
\textbf{Method} & \textbf{Selected role} & $\boldsymbol{TP/N}$ & \textbf{Recall} \\
\midrule
Generalist & -- & 35/50 & 0.70 \\
Fixed role & Domain & 38/50 & 0.76 \\
All roles & Generalist + Methods + Domain & 40/50 & 0.80 \\
\bottomrule
\end{tabular}
\caption{Working-only ParlaMint-GB Qwen3 8B detection check at balanced \(N=50\). The source is \texttt{parlamint\_gb\_fullsample\_eval100\_working\_v1}; CIs were not estimated for this interim check.}
\label{tab:working_parlamint_qwen_n50}
\end{table*}

\begin{table*}[t]
\centering
\scriptsize
\setlength{\tabcolsep}{3.2pt}
\renewcommand{\arraystretch}{1.05}
\begin{tabular}{@{}lcccc@{}}
\toprule
\textbf{Method} & \textbf{N=25 recall} & \textbf{N=50 recall} & \textbf{Added calls} & \textbf{Recall gain} \\
\midrule
Generalist & 0.76 & 0.70 & 25 & -0.06 \\
Fixed role & 0.76 & 0.76 & 25 & 0.00 \\
All roles & 0.80 & 0.80 & 75 & 0.00 \\
\bottomrule
\end{tabular}
\caption{Working-only sample-size efficiency check for ParlaMint-GB Qwen3 8B. In this interim run, recall does not improve from \(N=25\) to \(N=50\), while All roles requires three reviewer calls per packet.}
\label{tab:working_parlamint_efficiency}
\end{table*}
\fi

\subsection{RQ2 analysis set and operational accounting}

For each reviewer snapshot and corpus, expected calls equal eligible variants multiplied by three role prompts and three repetitions. Table~\ref{tab:rq2_accounting} records expected and completed calls, valid outputs, and cannot judge, abstention, and terminal failure counts. The reporting field below links failure types, tokens, latency, and cost to the same run identifiers.

\begin{table*}[t]
\centering
\scriptsize
\setlength{\tabcolsep}{3pt}
\renewcommand{\arraystretch}{1.08}
\begin{tabularx}{\textwidth}{@{}
>{\raggedright\arraybackslash}p{2.1cm}
>{\raggedright\arraybackslash}p{3.0cm}
>{\centering\arraybackslash}p{1.5cm}
>{\centering\arraybackslash}p{1.5cm}
>{\centering\arraybackslash}p{2.0cm}
>{\centering\arraybackslash}p{1.4cm}
>{\centering\arraybackslash}X@{}}
\toprule
\textbf{Corpus / analysis} & \textbf{Exact reviewer snapshot / class} &
\textbf{Source clusters} & \textbf{Eligible variants} &
\textbf{Expected / completed calls} & \textbf{Valid} &
\textbf{Cannot judge / abstain / terminal} \\
\midrule
Dreaddit test audit & \fillin{PRIMARY SNAPSHOT} / primary & \fillin{K} & \fillin{N} & \fillin{E/C} & \fillin{V} & \fillin{CJ/A/T} \\
Dreaddit test audit & \fillin{SENSITIVITY SNAPSHOT; ONE ROW EACH} & \fillin{K} & \fillin{N} & \fillin{E/C} & \fillin{V} & \fillin{CJ/A/T} \\
GoEmotions test audit & \fillin{PRIMARY SNAPSHOT} / primary & \fillin{K} & \fillin{N} & \fillin{E/C} & \fillin{V} & \fillin{CJ/A/T} \\
GoEmotions test audit & \fillin{SENSITIVITY SNAPSHOT; ONE ROW EACH} & \fillin{K} & \fillin{N} & \fillin{E/C} & \fillin{V} & \fillin{CJ/A/T} \\
CaChe held-out & \fillin{PRIMARY SNAPSHOT} / primary & \fillin{K} & \fillin{N} & \fillin{E/C} & \fillin{V} & \fillin{CJ/A/T} \\
CaChe held-out & \fillin{SENSITIVITY SNAPSHOT; ONE ROW EACH} & \fillin{K} & \fillin{N} & \fillin{E/C} & \fillin{V} & \fillin{CJ/A/T} \\
ParlaMint-GB held-out & \fillin{PRIMARY SNAPSHOT} / primary & \fillin{K} & \fillin{N} & \fillin{E/C} & \fillin{V} & \fillin{CJ/A/T} \\
ParlaMint-GB held-out & \fillin{SENSITIVITY SNAPSHOT; ONE ROW EACH} & \fillin{K} & \fillin{N} & \fillin{E/C} & \fillin{V} & \fillin{CJ/A/T} \\
\bottomrule
\end{tabularx}
\caption{RQ2 analysis and execution accounting.}
\label{tab:rq2_accounting}
\end{table*}

\fillin{REPORT ACTOR IDS, MODEL SNAPSHOTS, PROMPT AND SCHEMA HASHES, RUN DATES, OUTPUT SEALS, ANALYSIS EXPORT, FAILURE TYPES, INPUT/CACHED/OUTPUT/REASONING TOKENS, MEDIAN/IQR/TAIL LATENCY, AND COST FOR EACH ROW OF TABLE~\ref{tab:rq2_accounting}}.

\subsection{RQ2 role-by-flaw estimates}

Table~\ref{tab:rq2_role_flaw} organizes the recall quantities and descriptive stratification defined in Sections~\ref{sec:evaluation_methods} and~\ref{sec:evaluation_analysis} by error type for the primary reviewer snapshot. Each frozen sensitivity snapshot repeats the complete block.

\begin{table*}[t]
\centering
\tiny
\setlength{\tabcolsep}{2.2pt}
\renewcommand{\arraystretch}{1.12}
\begin{tabular}{@{}llcccccc@{}}
\toprule
\textbf{Corpus} & \textbf{Reviewer / method} &
\textbf{Unsupported evidence} & \textbf{Source concentration} &
\textbf{Counterevidence loss} & \textbf{Contextual flattening} &
\textbf{Unsupported abstraction} & \textbf{Overall} \\
\midrule
Dreaddit test & Researcher prompt & \fillin{TP/N [CI]} & \fillin{TP/N [CI]} & \fillin{TP/N [CI]} & \fillin{TP/N [CI]} & \fillin{TP/N [CI]} & \fillin{TP/N [CI]} \\
& Methods prompt & \fillin{TP/N [CI]} & \fillin{TP/N [CI]} & \fillin{TP/N [CI]} & \fillin{TP/N [CI]} & \fillin{TP/N [CI]} & \fillin{TP/N [CI]} \\
& Domain prompt & \fillin{TP/N [CI]} & \fillin{TP/N [CI]} & \fillin{TP/N [CI]} & \fillin{TP/N [CI]} & \fillin{TP/N [CI]} & \fillin{TP/N [CI]} \\
& WarrantRoute & \fillin{TP/N [CI]} & \fillin{TP/N [CI]} & \fillin{TP/N [CI]} & \fillin{TP/N [CI]} & \fillin{TP/N [CI]} & \fillin{TP/N [CI]} \\
& Fixed role (\fillin{ROLE}) & \fillin{TP/N [CI]} & \fillin{TP/N [CI]} & \fillin{TP/N [CI]} & \fillin{TP/N [CI]} & \fillin{TP/N [CI]} & \fillin{TP/N [CI]} \\
\addlinespace[1pt]
GoEmotions test & Researcher prompt & \fillin{TP/N [CI]} & \fillin{TP/N [CI]} & \fillin{TP/N [CI]} & \fillin{TP/N [CI]} & \fillin{TP/N [CI]} & \fillin{TP/N [CI]} \\
& Methods prompt & \fillin{TP/N [CI]} & \fillin{TP/N [CI]} & \fillin{TP/N [CI]} & \fillin{TP/N [CI]} & \fillin{TP/N [CI]} & \fillin{TP/N [CI]} \\
& Domain prompt & \fillin{TP/N [CI]} & \fillin{TP/N [CI]} & \fillin{TP/N [CI]} & \fillin{TP/N [CI]} & \fillin{TP/N [CI]} & \fillin{TP/N [CI]} \\
& WarrantRoute & \fillin{TP/N [CI]} & \fillin{TP/N [CI]} & \fillin{TP/N [CI]} & \fillin{TP/N [CI]} & \fillin{TP/N [CI]} & \fillin{TP/N [CI]} \\
& Fixed role (\fillin{ROLE}) & \fillin{TP/N [CI]} & \fillin{TP/N [CI]} & \fillin{TP/N [CI]} & \fillin{TP/N [CI]} & \fillin{TP/N [CI]} & \fillin{TP/N [CI]} \\
\addlinespace[1pt]
CaChe & Researcher prompt & \fillin{TP/N [CI]} & \fillin{TP/N [CI]} & \fillin{TP/N [CI]} & \fillin{TP/N [CI]} & \fillin{TP/N [CI]} & \fillin{TP/N [CI]} \\
& Methods prompt & \fillin{TP/N [CI]} & \fillin{TP/N [CI]} & \fillin{TP/N [CI]} & \fillin{TP/N [CI]} & \fillin{TP/N [CI]} & \fillin{TP/N [CI]} \\
& Domain prompt & \fillin{TP/N [CI]} & \fillin{TP/N [CI]} & \fillin{TP/N [CI]} & \fillin{TP/N [CI]} & \fillin{TP/N [CI]} & \fillin{TP/N [CI]} \\
& WarrantRoute & \fillin{TP/N [CI]} & \fillin{TP/N [CI]} & \fillin{TP/N [CI]} & \fillin{TP/N [CI]} & \fillin{TP/N [CI]} & \fillin{TP/N [CI]} \\
& Fixed role (\fillin{ROLE}) & \fillin{TP/N [CI]} & \fillin{TP/N [CI]} & \fillin{TP/N [CI]} & \fillin{TP/N [CI]} & \fillin{TP/N [CI]} & \fillin{TP/N [CI]} \\
\addlinespace[1pt]
ParlaMint-GB & Researcher prompt & \fillin{TP/N [CI]} & \fillin{TP/N [CI]} & \fillin{TP/N [CI]} & \fillin{TP/N [CI]} & \fillin{TP/N [CI]} & \fillin{TP/N [CI]} \\
& Methods prompt & \fillin{TP/N [CI]} & \fillin{TP/N [CI]} & \fillin{TP/N [CI]} & \fillin{TP/N [CI]} & \fillin{TP/N [CI]} & \fillin{TP/N [CI]} \\
& Domain prompt & \fillin{TP/N [CI]} & \fillin{TP/N [CI]} & \fillin{TP/N [CI]} & \fillin{TP/N [CI]} & \fillin{TP/N [CI]} & \fillin{TP/N [CI]} \\
& WarrantRoute & \fillin{TP/N [CI]} & \fillin{TP/N [CI]} & \fillin{TP/N [CI]} & \fillin{TP/N [CI]} & \fillin{TP/N [CI]} & \fillin{TP/N [CI]} \\
& Fixed role (\fillin{ROLE}) & \fillin{TP/N [CI]} & \fillin{TP/N [CI]} & \fillin{TP/N [CI]} & \fillin{TP/N [CI]} & \fillin{TP/N [CI]} & \fillin{TP/N [CI]} \\
\bottomrule
\end{tabular}
\caption{RQ2 recall by reviewer role and error type.}
\label{tab:rq2_role_flaw}
\end{table*}

\paragraph{Model and repetition sensitivity.}
\fillin{REPORT EVERY PRESPECIFIED MODEL AND REPETITION SENSITIVITY QUANTITY DEFINED IN SECTION~\ref{sec:evaluation_analysis} AND THE COMPLETE TABLE~\ref{tab:rq2_role_flaw} BLOCK}. These analyses are descriptive unless a later prospective amendment defines a closed confirmatory family and multiplicity correction.

Supplementary RQ1 reporting records missingness, failures, timing, robustness checks, reliability of the human reference labels, and both human reference conditions without adding another performance metric. Every estimate in this section must come from its corresponding frozen real-corpus export.

\subsection{RQ3 repair reporting}

RQ3 reporting gives the eligible and allocated packet counts, completed revisions, valid panel assessments, target-repair counts, accurate-material preservation, each collateral-error type, terminal failures, and successful repair without collateral error for every frozen condition. \fillin{REPORT THE RQ3 ACTOR, REVISION, FEEDBACK, PANEL, RUN, OUTPUT-SEAL, DENOMINATOR, EFFECT, INTERVAL, FAILURE, TOKEN, LATENCY, COST, AND ANALYSIS EXPORT IDENTIFIERS AND QUANTITIES}.
